% Vietnamese translation for OI Public Library
% Source-File: IOI中国国家候选队论文/2006/周戈林/周戈林.ppt
\documentclass[11pt]{article}
\usepackage{oipl-vn}

\title{Tư tưởng tương tự\\{\large Ghi chú theo slide}}
\author{Zhou Gelin}
\date{2006}

\begin{document}
\maketitle

\origtitle{Analogy thinking in informatics}
\sourcepath{IOI China candidate team papers / 2006 / Zhou Gelin / Zhou Gelin.ppt}

\section{Phạm vi}

Bộ trình chiếu này đi cùng bài viết về cách dùng phép tương tự để chuyển bài
toán lạ thành bài toán quen. Ghi chú dựa trên bản bài viết đã dịch và giữ ba
hướng chính của mạch trình bày.

\section{Tương tự có kiểm chứng}

Tương tự không phải là đoán mò theo bề ngoài. Ta cần xác định hai đối tượng
giống nhau ở quan hệ bản chất nào, rồi mới chuyển kết luận từ đối tượng cũ
sang đối tượng mới. Trong thuật toán, một phép tương tự tốt thường dẫn tới mô
hình, trạng thái hoặc chứng minh đúng đắn.

\section{Từ sự vật tới mô hình}

Một bài toán thực tế có thể giống một mô hình trừu tượng nếu ta chọn đúng đặc
trưng cần giữ lại. Ví dụ trong bài viết, địa hình được nhìn thành cấu trúc
cây theo các tầng độ cao; khi đó bài toán giảm số đỉnh núi trở thành một bài
quy hoạch động trên cây thay vì xử lý trực tiếp trên dãy độ cao.

\section{Từ thuật toán này sang thuật toán khác}

Nhiều thuật toán có cấu trúc gần nhau: luồng và ghép, cây và đoạn liên tiếp,
hình học và biểu thức số. Khi nhận ra quan hệ tương tự, ta có thể mượn cách
chứng minh, cách lưu trạng thái hoặc cách tối ưu từ bài quen thuộc. Tuy vậy,
mỗi phép chuyển đều cần kiểm tra lại các giả thiết, vì chỉ một điều kiện nhỏ
khác đi cũng có thể làm thuật toán cũ mất hiệu lực.

\section{Ghi nhớ}

Phép tương tự hữu ích nhất khi nó làm rõ cấu trúc ẩn. Sau khi tìm được đối
tượng quen thuộc, bước tiếp theo luôn là chứng minh rằng các thao tác trong mô
hình mới bảo toàn đúng yêu cầu của đề.

\end{document}
